% Options for packages loaded elsewhere
\PassOptionsToPackage{unicode}{hyperref}
\PassOptionsToPackage{hyphens}{url}
\documentclass[
]{article}
\usepackage{xcolor}
\usepackage[margin=1in]{geometry}
\usepackage{amsmath,amssymb}
\setcounter{secnumdepth}{-\maxdimen} % remove section numbering
\usepackage{iftex}
\ifPDFTeX
  \usepackage[T1]{fontenc}
  \usepackage[utf8]{inputenc}
  \usepackage{textcomp} % provide euro and other symbols
\else % if luatex or xetex
  \usepackage{unicode-math} % this also loads fontspec
  \defaultfontfeatures{Scale=MatchLowercase}
  \defaultfontfeatures[\rmfamily]{Ligatures=TeX,Scale=1}
\fi
\usepackage{lmodern}
\ifPDFTeX\else
  % xetex/luatex font selection
\fi
% Use upquote if available, for straight quotes in verbatim environments
\IfFileExists{upquote.sty}{\usepackage{upquote}}{}
\IfFileExists{microtype.sty}{% use microtype if available
  \usepackage[]{microtype}
  \UseMicrotypeSet[protrusion]{basicmath} % disable protrusion for tt fonts
}{}
\makeatletter
\@ifundefined{KOMAClassName}{% if non-KOMA class
  \IfFileExists{parskip.sty}{%
    \usepackage{parskip}
  }{% else
    \setlength{\parindent}{0pt}
    \setlength{\parskip}{6pt plus 2pt minus 1pt}}
}{% if KOMA class
  \KOMAoptions{parskip=half}}
\makeatother
\usepackage{graphicx}
\makeatletter
\newsavebox\pandoc@box
\newcommand*\pandocbounded[1]{% scales image to fit in text height/width
  \sbox\pandoc@box{#1}%
  \Gscale@div\@tempa{\textheight}{\dimexpr\ht\pandoc@box+\dp\pandoc@box\relax}%
  \Gscale@div\@tempb{\linewidth}{\wd\pandoc@box}%
  \ifdim\@tempb\p@<\@tempa\p@\let\@tempa\@tempb\fi% select the smaller of both
  \ifdim\@tempa\p@<\p@\scalebox{\@tempa}{\usebox\pandoc@box}%
  \else\usebox{\pandoc@box}%
  \fi%
}
% Set default figure placement to htbp
\def\fps@figure{htbp}
\makeatother
\setlength{\emergencystretch}{3em} % prevent overfull lines
\providecommand{\tightlist}{%
  \setlength{\itemsep}{0pt}\setlength{\parskip}{0pt}}
\usepackage{bookmark}
\IfFileExists{xurl.sty}{\usepackage{xurl}}{} % add URL line breaks if available
\urlstyle{same}
\hypersetup{
  pdftitle={Impact of Social Media Usage on Academic Performance of College Students},
  hidelinks,
  pdfcreator={LaTeX via pandoc}}

\title{Impact of Social Media Usage on Academic Performance of College
Students}
\author{}
\date{\vspace{-2.5em}2026-05-30}

\begin{document}
\maketitle

\section{Project Overview}\label{project-overview}

This project analyzes the relationship between social media usage and
academic performance among college students.

A dataset of 50 students was created containing:

\begin{itemize}
\tightlist
\item
  Social Media Hours
\item
  Study Hours
\item
  Sleep Hours
\item
  Attendance
\item
  Marks
\end{itemize}

The main objective is to determine whether increased social media usage
affects academic performance.

\section{Descriptive Statistics}\label{descriptive-statistics}

\begin{itemize}
\tightlist
\item
  Total Students: 50
\item
  Mean Social Media Usage: 4.28 hours/day
\item
  Median Social Media Usage: 4 hours/day
\item
  Maximum Usage: 8 hours/day
\item
  Minimum Usage: 1 hour/day
\end{itemize}

\section{Visual Analysis}\label{visual-analysis}

The histogram shows the distribution of social media usage among
students.

The boxplot summarizes the spread of social media usage and indicates
there are no significant outliers.

\section{Correlation Analysis}\label{correlation-analysis}

The calculated correlation between Social Media Hours and Marks is:

\textbf{-0.99}

This indicates a very strong negative relationship.

\section{Key Observation}\label{key-observation}

Students who spend more time on social media tend to score lower marks.

Students using 1--2 hours of social media generally scored above 85\%.

Students using 7--8 hours generally scored below 65\%.

\section{Conclusion}\label{conclusion}

The analysis suggests that excessive social media usage may negatively
impact academic performance.

A strong inverse relationship was observed between social media usage
and marks.

\end{document}
